% GENERATED FILE. Edit canonical lessons, then run npm run generate:books.
% canonical-source-hash: fnv1a64:5bbc091579785c26
% canonical-lessons: SA-C182-vismayah, SA-C182-sokah, SA-C182-santosah, SA-C182-krtajnata, SA-C182-smrtih

\chapter{Wonder, Grief, Contentment, Gratitude, Memory}
\label{ch:sa-wonder-grief-contentment-gratitude-memory}
\hlchaptermodality{182}

\begin{chapteropening}
\textbf{By the end of this chapter:} \emph{I can say wonder, grief, contentment, gratitude and memory.}

Complete the last lesson of chapter 182: I can say wonder, grief, contentment, gratitude and memory.
\end{chapteropening}

\section[\texorpdfstring{\sk{विस्मयः} (vismayaḥ)}{विस्मयः (vismayaḥ)}]{\sk{विस्मयः} (vismayaḥ) — wonder, surprise}

\label{lesson:SA-C182-vismayah}



\begin{quote}
Before the new one: say the Sanskrit for the pulse, then the Sanskrit for fitness, health.
\end{quote}

\subsection*{You'll want to know: \sk{विस्मयः}}
\textbf{\sk{विस्मयः}} — \emph{vismayaḥ} — \textquotedblleft{}wonder, surprise\textquotedblright{}.

\begin{grammarlens}[title={What you've built}]
One more word for this chapter.
\end{grammarlens}

\subsection*{Your turn}
\begin{itemize}
\raggedright
  \item \emph{Say it:} \emph{vismayaḥ}
  \item \emph{Say it:} \emph{vismayaḥ}, once more
  \item \emph{Say it:} say \emph{vismayaḥ}
  \item \emph{Recall:} say the Sanskrit for the pulse, then the Sanskrit for fitness, health, then say \emph{vismayaḥ} again
\end{itemize}

\begin{tcolorbox}[breakable,colback=teal!4,colframe=teal!35!black,title={Before you move on}]
What does \sk{विस्मयः} mean? (\textquotedblleft{}Wonder, surprise\textquotedblright{}.) Say it once more.
\end{tcolorbox}
\section[\texorpdfstring{\sk{शोकः} (śokaḥ)}{शोकः (śokaḥ)}]{\sk{शोकः} (śokaḥ) — grief}

\label{lesson:SA-C182-sokah}



\begin{quote}
Before the new one: say the Sanskrit for fitness, health, then the Sanskrit for wonder.
\end{quote}

\subsection*{You'll want to know: \sk{शोकः}}
\textbf{\sk{शोकः}} — \emph{śokaḥ} — \textquotedblleft{}grief\textquotedblright{}.

\begin{grammarlens}[title={What you've built}]
One more word for this chapter.
\end{grammarlens}

\subsection*{Your turn}
\begin{itemize}
\raggedright
  \item \emph{Say it:} \emph{śokaḥ}
  \item \emph{Say it:} \emph{śokaḥ}, once more
  \item \emph{Say it:} say \emph{śokaḥ}
  \item \emph{Recall:} say the Sanskrit for fitness, health, then the Sanskrit for wonder, then say \emph{śokaḥ} again
\end{itemize}

\begin{tcolorbox}[breakable,colback=teal!4,colframe=teal!35!black,title={Before you move on}]
What does \sk{शोकः} mean? (\textquotedblleft{}Grief\textquotedblright{}.) Say it once more.
\end{tcolorbox}
\section[\texorpdfstring{\sk{सन्तोषः} (santoṣaḥ)}{सन्तोषः (santoṣaḥ)}]{\sk{सन्तोषः} (santoṣaḥ) — contentment}

\label{lesson:SA-C182-santosah}



\begin{quote}
Before the new one: say the Sanskrit for wonder, then the Sanskrit for grief.
\end{quote}

\subsection*{You'll want to know: \sk{सन्तोषः}}
\textbf{\sk{सन्तोषः}} — \emph{santoṣaḥ} — \textquotedblleft{}contentment\textquotedblright{}.

\begin{grammarlens}[title={What you've built}]
One more word for this chapter.
\end{grammarlens}

\subsection*{Your turn}
\begin{itemize}
\raggedright
  \item \emph{Say it:} \emph{santoṣaḥ}
  \item \emph{Say it:} \emph{santoṣaḥ}, once more
  \item \emph{Say it:} say \emph{santoṣaḥ}
  \item \emph{Recall:} say the Sanskrit for wonder, then the Sanskrit for grief, then say \emph{santoṣaḥ} again
\end{itemize}

\begin{tcolorbox}[breakable,colback=teal!4,colframe=teal!35!black,title={Before you move on}]
What does \sk{सन्तोषः} mean? (\textquotedblleft{}Contentment\textquotedblright{}.) Say it once more.
\end{tcolorbox}
\section[\texorpdfstring{\sk{कृतज्ञता} (kṛtajñatā)}{कृतज्ञता (kṛtajñatā)}]{\sk{कृतज्ञता} (kṛtajñatā) — gratitude}

\label{lesson:SA-C182-krtajnata}



\begin{quote}
Before the new one: say the Sanskrit for grief, then the Sanskrit for contentment.
\end{quote}

\subsection*{You'll want to know: \sk{कृतज्ञता}}
\textbf{\sk{कृतज्ञता}} — \emph{kṛtajñatā} — \textquotedblleft{}gratitude\textquotedblright{}.

\begin{grammarlens}[title={What you've built}]
One more word for this chapter.
\end{grammarlens}

\subsection*{Your turn}
\begin{itemize}
\raggedright
  \item \emph{Say it:} \emph{kṛtajñatā}
  \item \emph{Say it:} \emph{kṛtajñatā}, once more
  \item \emph{Say it:} say \emph{kṛtajñatā}
  \item \emph{Recall:} say the Sanskrit for grief, then the Sanskrit for contentment, then say \emph{kṛtajñatā} again
\end{itemize}

\begin{tcolorbox}[breakable,colback=teal!4,colframe=teal!35!black,title={Before you move on}]
What does \sk{कृतज्ञता} mean? (\textquotedblleft{}Gratitude\textquotedblright{}.) Say it once more.
\end{tcolorbox}
\section[\texorpdfstring{\sk{स्मृतिः} (smṛtiḥ)}{स्मृतिः (smṛtiḥ)}]{\sk{स्मृतिः} (smṛtiḥ) — memory}

\label{lesson:SA-C182-smrtih}



\begin{quote}
Before the new one: say the Sanskrit for contentment, then the Sanskrit for gratitude.
\end{quote}

\subsection*{You'll want to know: \sk{स्मृतिः}}
\textbf{\sk{स्मृतिः}} — \emph{smṛtiḥ} — \textquotedblleft{}memory\textquotedblright{}.

\begin{grammarlens}[title={What you've built}]
One more word for this chapter.
\end{grammarlens}

\subsection*{Your turn}
\begin{itemize}
\raggedright
  \item \emph{Say it:} \emph{smṛtiḥ}
  \item \emph{Say it:} \emph{smṛtiḥ}, once more
  \item \emph{Say it:} say \emph{smṛtiḥ}
  \item \emph{Recall:} say the Sanskrit for contentment, then the Sanskrit for gratitude, then say \emph{smṛtiḥ} again
\end{itemize}

\begin{tcolorbox}[breakable,colback=teal!4,colframe=teal!35!black,title={Before you move on}]
What does \sk{स्मृतिः} mean? (\textquotedblleft{}Memory\textquotedblright{}.) Say it once more.
\end{tcolorbox}
